% Section 11
\section{Appendix C: Observer-Relativity and the Multi-Observer Problem}
The construction is naturally multi-observer rather than 
single-frame. Standard QFT in the S-matrix formulation 
assumes a single asymptotic observer who prepares the 
initial state and measures the final state; in the 
language of Section~\ref{sec:types}, there is one lab 
and one measurement at the end. Horizon physics 
requires a different framework: four observers 
$\mathcal{P}$, $\mathcal{R}$, $\mathcal{L}$, 
$\mathcal{F}$ each hold partial, causally 
complementary information about the same physical 
process. Neither $\mathcal{R}$ nor $\mathcal{L}$ alone 
can reconstruct the emission at $\mathcal{P}$; their 
two classical bits together can. The horizon is not 
a barrier in a single-observer description but a 
surface that makes the multi-observer structure of 
the construction sharp: it is the boundary at which 
$\mathcal{R}$'s and $\mathcal{L}$'s causal domains 
separate. This is not a limitation of the framework 
but its natural setting, and it is structurally the 
same situation that arises in Wigner's friend 
scenarios~\cite{brukner2018no} applied to 
gravitational systems, a point developed in the following section.

The multi-observer structure exposed here is not 
merely an analogy with Wigner's friend: it is 
precisely that scenario, instantiated in a 
gravitational context with causal separation enforced 
by a horizon rather than by a laboratory door.

A complete formalism for multi-observer quantum 
mechanics, one that specifies consistently how 
causally separated observers' measurement records 
combine into a single description, does not yet 
exist. The Wigner's friend no-go 
theorems~\cite{brukner2018no} and the 
Frauchiger--Renner result~\cite{frauchiger2018quantum} 
establish that the problem is sharp; the relational, 
QBist, and Everettian programs each dissolve it 
differently rather than solving it. The present 
construction operates within a single observer's 
description but does not treat the gap as a 
liability: the observer-independence of the 
Bekenstein--Hawking area law emerges here not from 
a completed multi-observer formalism but from the 
common causal ancestry of $\mathcal{P}$, which 
enforces intersubjective consistency between 
$\mathcal{R}$'s and $\mathcal{L}$'s records without 
requiring an objective wavefunction. The 
Frauchiger--Renner and Wigner's friend problems 
remain open in general; the construction shows that 
causal ancestry is sufficient to ground 
intersubjective agreement in this specific setting.

Whether the relational~\cite{rovelli1996relational}, 
QBist~\cite{fuchs2014introducing}, or Everettian 
frameworks are inequivalent in this setting is an 
open question with physical consequences:
different resolutions imply different answers to whether 
$\mathcal{R}$'s microstate count and $\mathcal{L}$'s 
microstate count are consistent as objective facts 
or only as relational ones. The information-theoretic 
objects introduced here, the spatial bit, the PVM 
projection, the handoff event, the matter-memory, are 
the vocabulary in which that question must eventually 
be posed, but they do not answer it. Multi-observer 
quantum mechanics in the presence of horizons is 
largely uncharted, and the present construction 
should be read partly as a map of that problem space.
